\section{Day 14: Independence of RVs, Expected Value (Nov 03, 2025)}

\subsection{Independence of Random Variables}

\begin{definition}
For general two random variables $X$ and $Y$, they are independent if the events $\{ X \in B \}$ and $\{ Y \in C \}$ are independent for all subsets $B, C \in \mathbb{R}$, meaning
\[
    P(X \in B, Y \in C) = P(X \in B)P(Y \in C)
\]
\begin{itemize}
\item    If $X, Y$ are discrete, given $p_{X, Y}(x, y)$, we can define the conditional distribution
\[
    p_{Y \mid X}(y \mid x) = \frac{p_{X, Y}(x, y)}{p_X(x)}
\]
\item    If $X, Y$ are abs. continuous, given $f_{X, Y}$ (joint density function for $X, Y$), define
\[
    f_{Y \mid X}(y \mid x) = \frac{f_{X, Y}(x, y)}{f_X(x)}
\]
called the conditional probability density of $Y$ given the occurrence of the value $x$ of $X$. 
\end{itemize}
\end{definition}

\begin{simplethm}
    For absolutely continuous random variables $X, Y$, $X$ and $Y$ are independent iff $f_{X, Y}(x, y) = f_X(x)f_Y(y)$ for all $x, y \in \mathbb{R}$. 
\end{simplethm}

\begin{simplethm}
    Suppose $X, Y$ are discrete and independent. Given $p_{X, Y}(x, y)$, and $p_X(x) > 0$, then 
    \[
        p_{Y \mid X}(y \mid x) = p_Y(y)
    \]
\end{simplethm}

\begin{simplethm}
    Suppose $X, Y$ are continuous and independent. If $f_{X, Y}(x, y)$, $p_X(x) > 0$, then 
    \[
        p_{Y \mid X}(y \mid x) = p_Y(y)
    \]
\end{simplethm}

\noindent Chapter 2 of the textbook ends here.

\subsection{Expected Value}

\begin{definition}[Expected Value]
    If $X$ is a discrete random variable, the expected value $E(X)$\footnote{strictly speaking, in this course we use $\mathrm{E}(X)$, so mathrm should be used on E} is given by
    \[
    E(X) = \sum_{x \in \mathbb{R}} x \, p_X(x)
    \]
\end{definition}

\noindent In other courses/textbooks, the alternative notation $\mathbb{E}[X], \mathbb{E}(X)$ may be used. 
